\section{结论与应用建议}
多能源转型的安全性由建设兑现、运行交付与资源准备共同决定。三问形成可逐项核查的决策链。

现状上，2010—2024年一次能源消费增长65.3\%，化石服务与新能源扩张并行；主体模型通过独立料源、建设队列与交付约束承接这一矛盾。双碳推演中，基准至2060窗口末端仍未确认达峰，政策强化将2060剩余净排由12.639降至9.325 Gt；主体主动反馈的终点减排伴随跨期服务代价，评价需同时观察排放与交付。

投资上，折中路径比最低观察可靠资本多投入13.663万亿元、减少2060净排7.520 Gt，但原48组联合扰动仅3组年度可靠。保供修复将其增至20/48，另48组独立设计中21组达标，名义情景的折现资本增加26.526万亿元，2060年净排增加2.037 Gt。进口—工期网格定位不晚于2030年投产的油项目台阶；年度方案下钻小时层后，方案Ⅰ、方案Ⅲ在主窗满足电力交付，方案Ⅱ仍有8.827 TWh缺供。196条建设窗口对照中，窗口内扩建改善6组2030油缺口，却未提高20/48完整年度可靠数；全期扩大油气建设降至19/48。应先核对截止前实际可投产的队列，再配置国内建设与外部资源。

运行上，八组累计缺供均达到对应离线下界；原相位四段夜间守恒界为5.170 TWh。增加功率可改变小时缺口，改变累计电量还需可用供给、可交付储能或服务需求的变化。24臂联合约束进一步表明，八组最低费用均可同时守住原峰值；排放约束使降费空间从12.861—20.924亿元收缩为六组698.316—2738.648万元、两组数值等费。按累计能量、峰值、排放、费用与末态共同选择排程。

\begin{table}[H]\centering\caption{按瓶颈选择动作及其已检验作用}
\begin{tabularx}{\linewidth}{p{.15\linewidth}p{.23\linewidth}X}\toprule
识别的瓶颈 & 对应动作 & 完整路径中的实际作用\\\midrule
建设兑现与料源 & 分期资本及技术/供应订单 & 年度门槛、进口—工期网格与小时交付共同筛查所选组合\\
时序分配 & 联合协调储能、发电与偿还 & 对可达全电源余量局部规则，共同满电终点下原功率主窗少缺供109.455 GWh；32格中21格进一步改善\\
交付功率 & 功率增加50\%，能量固定 & 32格中29格削峰、5格减少累计缺供；原48小时设计削峰78.634 GW\\
跨事件准备 & 按时滞、渠道和资源目标补库 & 原付费合同增购末库存，新增现金、资产与溢价分账；四方案对原无事件目标的第二窗口库存仍低于95\%\\
\bottomrule\end{tabularx}\end{table}
三类运行动作分别对应时序、交付上限与资源补充。更强局部基线、事件信息、完整事件网格和2.3小时电池挑战给出收益的具体条件；库存以最近事件前、固定初始和同小时无事件三种参照共同解释。实际原生运行、完整跨实现路径与独立守恒账使方案、动作和结果衔接。工程应用时，以地区负荷、真实到货、库存和项目工期替换相应聚合输入，再按同一账本复验候选。

\subsection{五类协同对策与实施顺序}
\textbf{能源结构优化：先核投产窗口，再配置资本。}按截至目标年可兑现能力审查油气与电源订单。油工期为4、5、6年时，不晚于2030年投产的订单截止分别为2026、2025、2024年；窗口外的即时缺口应由已有队列、渠道及可用库存共同保障。P2改善6组单年油缺口而未提高全期可靠数，P3又出现新增电力失守，说明扩建须和其他资产交付共同评价。

\textbf{多能协同调度：同时检查能量、峰值与末态。}以实际状态滚动协调发电、储能和偿还，共同满电终点下相对强局部规则仍少缺供109.455 GWh。功率增加50\%在29/32格削峰、仅5格减少累计缺供；原相位已到四夜能量下界，应把功率配置与可交付能量分开决策，并直接加入峰值和排放约束选择费用排程。

\textbf{储备与应急建设：按到货时滞和资源目标补库。}把采购现金、末库存资产及溢价分别列账，本组1666.523亿元现金形成1111.016亿元资产和555.508亿元资源净增费。储备目标同时核对最近事件前、固定初始与同小时无事件库存；按跨境通道和合同到货条件落实资源，持续检查下一事件前的准备。

\textbf{市场机制设计：将稀缺与碳信号接到交付账。}联合评价价格规则诱发的订单、投产和需求变化。基准阶跃档位由800升至1200元/MWh时，2060年能源CO$_2$少0.200 Gt，累计未供服务却由19.608增至35.427 PWh；同档平滑规则为20.687 PWh。设计价格与碳反馈时，应同步比较供能、工程资本及资源费用，避免以终点减排替代完整服务评价。

\textbf{科技创新：以能量和工业净排缺口确定研发方向。}同一功率下，电池由4小时缩至2.3小时使原事件累计缺供增加396.654 GWh，可交付电量是独立的配置对象。基准工业过程与固定碳汇假设下仍有175.346 Mt结构净排下界，需要将工业深度减排或新增移除作为后续技术选项单独评价，不能只靠现有发电结构调整兑现净零。

实施上，近期先使用现有运行数据重排调度、明确库存目标和有效订单窗口；中期落实渠道、合同与已批准建设队列；长期依据地区负荷、到货和技术成本评价新增能量容量及工业减排。每一步均以实际交付、完整资源账和资本代价复核。
